\documentclass[9pt,a4paper]{article}
\usepackage{iftex}

\ifPDFTeX
  \usepackage[utf8]{inputenc}
  \IfFileExists{l7xenc.def}{%
    \usepackage[L7x,T1]{fontenc}%
  }{%
    \usepackage[T1]{fontenc}%
  }
  \IfFileExists{lithuanian.ldf}{%
    \usepackage[lithuanian]{babel}%
  }{}
  \usepackage{lmodern}
  % Atsarginis UTF-8 lietuviškų raidžių atvaizdavimas, jei kompiliuojama be L7x paketo
  \makeatletter
  \@ifundefined{DeclareUnicodeCharacter}{}{%
    \DeclareUnicodeCharacter{0104}{\k{A}}% Ą
    \DeclareUnicodeCharacter{0105}{\k{a}}% ą
    \DeclareUnicodeCharacter{010C}{\v{C}}% Č
    \DeclareUnicodeCharacter{010D}{\v{c}}% č
    \DeclareUnicodeCharacter{0118}{\k{E}}% Ę
    \DeclareUnicodeCharacter{0119}{\k{e}}% ę
    \DeclareUnicodeCharacter{0116}{\.{E}}% Ė
    \DeclareUnicodeCharacter{0117}{\.{e}}% ė
    \DeclareUnicodeCharacter{012E}{\k{I}}% Į
    \DeclareUnicodeCharacter{012F}{\k{i}}% į
    \DeclareUnicodeCharacter{0160}{\v{S}}% Š
    \DeclareUnicodeCharacter{0161}{\v{s}}% š
    \DeclareUnicodeCharacter{0172}{\k{U}}% Ų
    \DeclareUnicodeCharacter{0173}{\k{u}}% ų
    \DeclareUnicodeCharacter{016A}{\={U}}% Ū
    \DeclareUnicodeCharacter{016B}{\={u}}% ū
    \DeclareUnicodeCharacter{017D}{\v{Z}}% Ž
    \DeclareUnicodeCharacter{017E}{\v{z}}% ž
  }
  \makeatother
\else
  \usepackage{fontspec}
\fi

\usepackage[margin=1.2cm]{geometry}
\usepackage{amsmath,amssymb,amsfonts}
\usepackage{multicol}
\usepackage{booktabs}
\usepackage{array}
\usepackage{graphicx}
\usepackage{xcolor}
\usepackage{enumitem}
\usepackage{fancyhdr}

\setlist{nosep, leftmargin=*}
\pagestyle{fancy}
\fancyhf{}
\lhead{\textbf{Edukometrijos pagrindai: SEM, Kelių analizė ir CFA (S2 paruoštukė)}}
\rhead{VU MIF | Prof. V. Čekanavičius}
\cfoot{\thepage}

\setlength{\columnsep}{16pt}
\setlength{\columnseprule}{0.2pt}

% Tiksli dėžučių pločio kontrolė, apsauganti nuo stulpelio perpildymo (overfull hbox)
\newcommand{\boxalert}[2]{%
\noindent\fbox{%
\parbox{\dimexpr\linewidth-2\fboxsep-2\fboxrule\relax}{%
\textbf{\textcolor{red!80!black}{#1:}} #2%
}%
}\vspace{1.5mm}%
}

\newcommand{\boxtip}[2]{%
\noindent\fbox{%
\parbox{\dimexpr\linewidth-2\fboxsep-2\fboxrule\relax}{%
\textbf{\textcolor{blue!80!black}{#1:}} #2%
}%
}\vspace{1.5mm}%
}

\title{\vspace{-1.2cm}\Large\textbf{SEM, Kelių analizė ir CFA: Pilnas formulių konspektas}\\\large Skirta S2 atsiskaitymui ir teorijai (\textit{pagal „Prieš S2(5)“ ir 4 SEM skaidres})\vspace{-0.5cm}}
\author{Parengta pagal prof. V. Čekanavičiaus paskaitų medžiagą}
\date{}

\begin{document}
\maketitle
\thispagestyle{fancy}

\begin{multicols*}{2}

\section{Esmė ir modelių hierarchija}
\textbf{Struktūrinių lygčių modeliavimas} (angl. \textit{Structural Equation Modeling, SEM}) -- daugiamatės statistikos metodas, tikrinantis, ar tyrėjo iškeltas kintamųjų priklausomybių modelis yra suderintas su stebimais duomenimis.

\textbf{SEM apjungia 3 sritis}:
\begin{enumerate}
    \item \textbf{Tiesinę regresiją} (vieno endogeninio kintamojo priklausomybė nuo kelių egzogeninių).
    \item \textbf{Kelių analizę} (\textit{Path Analysis, PA}) -- priežastiniai ryšiai tarp kelių stebimų kintamųjų (kai kintamasis vienoje lygtyje yra priklausomas, kitoje -- nepriklausomas).
    \item \textbf{Patvirtinančiąją faktorinę analizę} (\textit{CFA}) -- ryšiai tarp stebimų indikatorių ir nematuojamų (\textbf{latentinių}) faktorių.
\end{enumerate}

\section{Kelių diagramos elementai}
\begin{itemize}
    \item \textbf{Stačiakampis} $\square$: \textbf{stebimas} kintamasis ($X, Y$).
    \item \textbf{Ovalas / apskritimas} $\bigcirc$: \textbf{latentinis} kintamasis ($F, \xi, \eta$) arba liekamoji paklaida ($e, \zeta, \delta$).
    \item \textbf{Vienpusė rodyklė} $\longrightarrow$: \textbf{priežastinis ryšys} (regresija). Rodyklė eina nuo priežasties link pasekmės ($X \to Y$).
    \item \textbf{Dvipusė lanko formos rodyklė} $\longleftrightarrow$: \textbf{kovariacija / koreliacija} (be krypties priežastingumo).
\end{itemize}

\subsection{Kintamųjų tipai}
\begin{itemize}
    \item \textbf{Egzogeninis kintamasis} (nepriklausomas): neturi jokių įeinančių vienpusių rodyklių (iš jo rodyklės tik išeina). Jo dispersija ir kovariacijos su kitais egzogeniniais kintamaisiais vertinamos tiesiogiai.
    \item \textbf{Endogeninis kintamasis} (priklausomas): turi bent vieną įeinančią vienpusę rodyklę $\to$. Kiekvienas endogeninis kintamasis \textbf{privalo turėti liekamąją paklaidą} ($e$), kurios dispersija $\text{Var}(e)$ yra vertinama.
\end{itemize}

\section{Modelio identifikacija ir laisvės laipsniai ($df$)}

\subsection{Žinomų parametrų skaičius}
Modeliui vertinti naudojama empirinė kovariacijų matrica $S$. Jei turime $k$ stebimų kintamųjų, matrica $S$ turi $k \times k$ elementų, tačiau dėl simetrijos nepriklausomos informacijos kiekis yra:
\begin{equation*}
    p^* = \frac{k(k + 1)}{2}
\end{equation*}

\subsection{Nežinomų parametrų skaičius ($q$)}
Laisvi vertinami parametrai modelio viduje:
\begin{itemize}
    \item Kelių koeficientai (regresijos svoriai).
    \item Egzogeninių kintamųjų dispersijos.
    \item Endogeninių kintamųjų liekanų dispersijos $\text{Var}(e)$.
    \item Kovariacijos tarp egzogeninių kintamųjų ar liekanų.
\end{itemize}

\subsection{Laisvės laipsniai ($df$)}
\begin{equation*}
    df = p^* - q = \frac{k(k + 1)}{2} - q
\end{equation*}

\subsection{Identifikacijos būsenos}
\begin{enumerate}
    \item \textbf{Poidentifikuotas} (\textit{Under-identified}): $df < 0$ ($p^* < q$). Lygčių mažiau nei nežinomųjų. Parametrų įvertinti neįmanoma!
    \item \textbf{Tiksliai identifikuotas / sotusis} (\textit{Just-identified / Saturated}): $df = 0$ ($p^* = q$). Parametrai randami vienareikšmiškai, bet $\chi^2 = 0$, $p = 1$ -- modelio tinkamumo patikrinti neįmanoma.
    \item \textbf{Viršidentifikuotas} (\textit{Over-identified}): $df > 0$ ($p^* > q$). Tai yra \textbf{būtina sąlyga} geram SEM tyrimui! Galima tikrinti modelio tinkamumą duomenims ($\chi^2$, RMSEA ir kt.).
\end{enumerate}

\subsection{Latentinių kintamųjų mastelio fiksavimas}
Kadangi latentinis kintamasis neturi fizinės skalės, jo identifikavimui privaloma atlikti vieną iš dviejų:
\begin{itemize}
    \item \textbf{Referentinis indikatorius} (\textit{Marker variable}): vieno pasirinkto indikatoriaus kelias fiksuojamas lygiu 1 ($\lambda_1 = 1$). Tai numatytasis nustatymas \texttt{lavaan}.
    \item \textbf{Faktoriaus standartizavimas}: faktoriaus dispersija fiksuojama lygi 1 ($\text{Var}(F) = 1$, komanda \texttt{std.lv = TRUE}).
\end{itemize}

\section{Wright'o kelių sekimo taisyklės ir poveikiai}

\subsection{Tiesioginis, netiesioginis ir bendrasis poveikis}
\begin{itemize}
    \item \textbf{Tiesioginis poveikis} (\textit{Direct effect}): koeficientas ant tiesioginės rodyklės tarp dviejų kintamųjų: $X \xrightarrow{\beta} Y \implies \text{Direct} = \beta$.
    \item \textbf{Netiesioginis poveikis} (\textit{Indirect effect}): poveikis per tarpinį kintamąjį (mediatorių):
    \begin{align*}
        &X \xrightarrow{a} M \xrightarrow{b} Y \\
        &\text{Indirect} = a \cdot b
    \end{align*}
    Jei yra keli netiesioginiai keliai, jų sandaugos sudedamos: $\sum a_j b_j$.
    \item \textbf{Bendrasis poveikis} (\textit{Total effect}):
    \begin{align*}
        \text{Total} &= \text{Direct} + \text{Indirect} \\
        &= \beta + \sum a_j b_j
    \end{align*}
\end{itemize}

\subsection{Wright'o kelių sekimo taisyklės (\textit{Tracing rules})}
Koreliacija tarp bet kurių dviejų kintamųjų standartizuotame modelyje randama sudedant visų leistinų jungiančių kelių koeficientų sandaugas.

\textbf{Leistino kelio taisyklės}:
\begin{enumerate}
    \item Negalima eiti per tą patį kintamąjį du kartus (nėra kilpų).
    \item Negalima eiti pirmyn rodykle, o po to atgal ($\to \leftarrow$ draudžiama!). Galima eiti atgal, o po to pirmyn ($\leftarrow \to$ leidžiama).
    \item Kelyje leidžiama daugiausiai viena dvipusė rodyklė $\leftrightarrow$ (kovariacija).
\end{enumerate}

\textbf{Pavyzdys iš paskaitos}:
Jei $X \xrightarrow{0.35} M \xrightarrow{0.67} Y$ ir $X \xrightarrow{0.45} Y$:
\begin{align*}
    \text{cor}(X, Y) &= \text{Direct} + \text{Indirect} \\
    &= 0.45 + 0.35 \times 0.67 = 0.6845
\end{align*}

\section{SEM matricinis pavidalas (LISREL ir RAM)}

\subsection{LISREL struktūra}
\begin{itemize}
    \item \textbf{Struktūrinis modelis} (ryšiai tarp latentinių dydžių):
    \begin{equation*}
        \eta = B \eta + \Gamma \xi + \zeta
    \end{equation*}
    kur $\eta$ -- endogeniniai latentiniai, $\xi$ -- egzogeniniai latentiniai, $B$ -- endogeninių ryšiai, $\Gamma$ -- egzogeninių poveikiai, $\zeta$ -- paklaidos.
    \item \textbf{Matavimo modeliai} (ryšiai su stebimais $x, y$):
    \begin{align*}
        y &= \Lambda_y \eta + \epsilon \\
        x &= \Lambda_x \xi + \delta
    \end{align*}
\end{itemize}

\subsection{RAM (Reticular Action Model) formulė}
Visi kintamieji sujungiami į vektorių $v$:
\begin{equation*}
    v = A v + u
\end{equation*}
kur $A$ -- asimetrinių kelių matrica ($\to$), $S$ -- simetrinių ryšių matrica (kovariacijos ir dispersijos $\leftrightarrow$), $F$ -- filtrinė matrica ($x = F v$).
Atkurta kovariacijų matrica:
\begin{equation*}
    \Sigma(\theta) = F (I - A)^{-1} S (I - A)^{-T} F^T
\end{equation*}

\section{Modelio tinkamumo rodikliai (\textit{Fit Measures})}

\subsection{1. Chi-kvadrato ($\chi^2$) testas}
Tikrina nulinę hipotezę $H_0: \Sigma(\theta) = \Sigma$ (atkurtas modelis idealiai atitinka populiacijos kovariacijų matricą).
\begin{itemize}
    \item \textbf{Siekiama reikšmė}: $p > 0.05$ (nėra reikšmingo skirtumo tarp modelio ir duomenų $\implies$ modelis tinka!).
    \item \textbf{Prof. Čekanavičiaus dėsnis dėl imties dydžio}:
    Jei imtis maža ($n < 200$), norime $p > 0.05$. Jei imtis didelė ($n > 400-500$), $\chi^2$ automatiškai tampa statistiškai reikšmingas ($p < 0.05$) net prie menkų nuokrypių. Tuomet $\chi^2$ $p$-reikšme nebepasitikima ir vertinama pagal indeksus!
\end{itemize}

\subsection{2. Normuotas Chi-kvadratas ($NC$)}
\begin{align*}
    NC &= \frac{\chi^2}{df} = \frac{\text{Test statistic}}{\text{df}}
\end{align*}
\textbf{Kriterijai}:
\begin{itemize}
    \item $NC < 2$ -- puikus modelio tinkamumas (labai gerai).
    \item $NC < 3$ -- priimtinas tinkamumas (iš bėdos tinka).
    \item $NC > 5$ -- prastas modelis.
\end{itemize}

\subsection{3. RMSEA (Root Mean Square Error of Approx.)}
\begin{equation*}
    RMSEA = \sqrt{\max\left(0,\, \frac{\chi^2 - df}{(n - 1)df}\right)}
\end{equation*}
\textbf{Interpretacija}:
\begin{itemize}
    \item $RMSEA < 0.05$ -- modelis \textbf{labai gerai} tinka duomenims.
    \item $0.05 \le RMSEA \le 0.08$ -- patenkinamas, priimtinas tinkamumas.
    \item $RMSEA > 0.08$ -- prastas tinkamumas ($> 0.10$ modelis netinkamas).
    \item Tikrinama $H_0: RMSEA \le 0.05$: norime $p > 0.05$.
\end{itemize}

\subsection{4. Lyginamieji indeksai (CFI, TLI, SRMR)}
\begin{itemize}
    \item \textbf{CFI (Comparative Fit Index)}:
    Palygina tiriamąjį modelį su nepriklausomu (nuliniu).
    \textbf{Geras}: $CFI > 0.95$; priimtinas: $CFI > 0.90$.
    \item \textbf{TLI (Tucker-Lewis Index)}:
    Baudžia už pernelyg sudėtingą modelį.
    \textbf{Geras}: $TLI > 0.95$; priimtinas: $TLI > 0.90$.
    \item \textbf{SRMR (Std. Root Mean Square Residual)}:
    Standartizuotas liekanų vidurkis.
    \textbf{Geras}: $SRMR < 0.08$.
    \item \textbf{AIC (Akaike Information Criterion)}:
    Naudojamas ne vieno modelio vertinimui, o kelių modelių \textbf{lyginimui}. Geresnis tas modelis, kurio \textbf{AIC yra mažesnis}.
\end{itemize}

\section{Parametrų įvertinimas ir interpretavimas}

\subsection{Reikšmingumo tikrinimas}
Kiekvienam vertinamam parametrui $\theta$ pateikiama:
\begin{itemize}
    \item Įvertis (\textit{Estimate}) $\hat{\theta}$.
    \item Standartinė paklaida (\textit{Std.Err}) $SE(\hat{\theta})$.
    \item $z$-statistika: $z = \frac{\hat{\theta}}{SE(\hat{\theta})}$.
    \item $p$-reikšmė ($P(>|z|)$): jei $p < 0.05$, parametras yra \textbf{statistiškai reikšmingas}.
\end{itemize}

\subsection{Standartizuoti vs nestandartizuoti parametrai}
\begin{itemize}
    \item \textbf{Nestandartizuoti} (\texttt{Estimate}): rodo pokytį originaliais kintamųjų vienetais. Naudojami struktūrinių lygčių užrašymui ir prognozavimui.
    \item \textbf{Standartizuoti} (\texttt{Std.all}): visi kintamieji transformuoti į $z$-skalę (vidurkis 0, dispersija 1).
    \textbf{Nauda}: leidžia tiesiogiai palyginti, kuris kintamasis daro \textbf{stipriausią santykinį poveikį}.
\end{itemize}

\boxalert{Dėstytojo pastabos dėl parametrų interpretacijos}{%
\begin{enumerate}
    \item \textbf{Nereikšmingos paklaidų dispersijos} ($p > 0.05$ prie $\text{Var}(e)$): tai nėra blogai! Tai rodo, kad priklausomybė sumodeliuota taip tiksliai, jog liekamoji dispersija praktiškai lygi 0.
    \item \textbf{Blogas modelis (Heywood atvejai)}: jei gaunama neigiama dispersija ($\text{Var}(e) < 0$) arba koreliacija $|r| > 1$, modelis yra blogai specifikuotas arba duomenys netinka.
    \item \textbf{Koreliacijos grafike}: rodo bendrą latentinių ar egzogeninių kintamųjų ryšį, nesusijusį su priežastine kryptimi.
\end{enumerate}
}

\section{R ir \texttt{lavaan} komandų paruoštukė}

\subsection{1. Modelio sintaksė}
\begin{center}
\resizebox{\columnwidth}{!}{%
\begin{tabular}{lll}
\toprule
\textbf{Operatorius} & \textbf{Reikšmė} & \textbf{Pavyzdys} \\
\midrule
\texttt{\~{}} & Regresija (kelias) & \texttt{Y \~{} X1 + X2} \\
\texttt{=\~{}} & Latentinis faktorius (CFA) & \texttt{F1 =\~{} A + B + C} \\
\texttt{\~{}\~{}} & Kovariacija / dispersija & \texttt{F1 \~{}\~{} F2} arba \texttt{X \~{}\~{} X} \\
\texttt{:=} & Naujas parametras & \texttt{ind := a*b} \\
\bottomrule
\end{tabular}%
}
\end{center}

\subsection{2. Modelio paleidimas su pradiniais duomenimis}
\begin{footnotesize}
\begin{verbatim}
library(lavaan)
library(lavaanPlot)
df <- na.omit(df)
mod <- '
  trstun ~ trstep
  happy ~ trstun + stfhlth
  stflife ~ trstun + happy + stfhlth
'
fit <- sem(mod, data=df, std.lv=TRUE)
\end{verbatim}
\end{footnotesize}

\subsection{3. Modelio paleidimas su kovariacijų matrica}
Kai duota kovariacijų matrica ir stebinių skaičius $n$:
\begin{footnotesize}
\begin{verbatim}
cov_mat <- matrix(c(
  7.0, 4.7, 4.4, 1.8, 1.1, 0.7,
  4.7, 7.3, 3.0, 2.0, 1.3, 0.8,
  4.4, 3.0, 5.2, 1.5, 0.9, 1.0,
  1.8, 2.0, 1.5, 4.7, 2.5, 1.5,
  1.1, 1.3, 0.9, 2.5, 3.0, 0.9,
  0.7, 0.8, 1.0, 1.5, 0.9, 5.5
), nrow=6, byrow=TRUE,
dimnames=list(
  c("A","B","C","D","E","F"),
  c("A","B","C","D","E","F")))

mod_cfa <- '
  F1 =~ A + B + C
  F2 =~ D + E + F
  F1 ~~ F2
'
fit_cfa <- sem(mod_cfa,
  sample.cov=cov_mat,
  sample.nobs=170)
\end{verbatim}
\end{footnotesize}

\subsection{4. Rezultatų išvedimas ir grafikas}
\begin{footnotesize}
\begin{verbatim}
# Išsami suvestinė su indeksais
summary(fit, standardized=TRUE,
        fit.measures=TRUE)

# Kelio diagrama su standartizuotais svoriais
lavaanPlot(model=fit, coefs=TRUE,
           stand=TRUE, covs=TRUE)
\end{verbatim}
\end{footnotesize}

\section{Praktiniai pavyzdžiai iš „Prieš S2(5)“}

\subsection{1 Pavyzdys: Kelių analizė (\textit{trust.sav})}
\textbf{Kintamieji}:
\begin{itemize}
    \item Egzogeniniai: $trstep$ (pasitikėjimas EP), $stfhlth$ (pasit. sveikatos apsauga).
    \item Endogeniniai: $trstun$ (pasit. JT), $happy$ (laimingumas), $stflife$ (pasitenkinimas gyvenimu).
\end{itemize}

\textbf{Struktūrinės lygtys su įverčiais}:
\begin{align*}
    trstun &= 0.635 \cdot trstep + e_1 \\
    happy &= 0.157 \cdot trstun + 0.156 \cdot stfhlth + e_2 \\
    stflife &= 0.165 \cdot trstun + 0.723 \cdot happy \\
    &\quad + 0.133 \cdot stfhlth + e_3
\end{align*}
Liekanų dispersijos: $\text{Var}(e_1)=3.947$, $\text{Var}(e_2)=2.692$, $\text{Var}(e_3)=2.338$.

\textbf{Tinkamumo vertinimas}:
\begin{itemize}
    \item $\chi^2 = 1.613$, $df = 3$, $p = 0.656 > 0.05 \implies$ modelis puikiai tinka duomenims!
    \item $NC = \frac{1.613}{3} = 0.538 < 2 \implies$ puiku.
    \item $RMSEA = 0.000 < 0.05 \implies$ tobulas.
\end{itemize}

\textbf{Netiesioginis poveikis}:
Poveikis nuo $trstun$ į $stflife$ per $happy$:
\begin{align*}
    \text{Indirect} &= 0.157 \times 0.723 = 0.114 \\
    \text{Total} &= 0.165 + 0.114 = 0.279
\end{align*}

\subsection{2 Pavyzdys: CFA su kovariacine matrica}
Duota 6 kintamųjų kovariacinė matrica ($k=6$, $n=170$).
\begin{itemize}
    \item Žinomos informacijos: $p^* = \frac{6 \times 7}{2} = 21$.
    \item Laisvi parametrai: 4 faktorių svoriai (2 fiksuoti į 1) + 6 paklaidų dispersijos + 2 faktorių dispersijos + 1 kovariacija $F_1 \leftrightarrow F_2 \implies q = 13$.
    \item Laisvės laipsniai: $df = 21 - 13 = 8$.
    \item $\chi^2 = 10.014$, $df = 8$, $p = 0.264 > 0.05 \implies$ modelis tinka duomenims.
    \item $RMSEA = 0.038 < 0.05 \implies$ geras modelis.
    \item Koreliacija tarp $F_1$ ir $F_2$ yra $0.375$ (vidutinė).
\end{itemize}

\section{Egzamino patarimai ir atsakymai}
\begin{enumerate}
    \item \textbf{Kada kintamieji endogeniniai?} Kai į juos įeina bent viena rodyklė $\to$. Jie visada turi paklaidą $e$.
    \item \textbf{Kaip atsakyti apie $\chi^2$?}
    Jei $p > 0.05$, rašome: „Modelis suderintas su duomenimis, duomenys modeliui neprieštarauja“. Jei $p < 0.05$, bet $n$ didelė, žiūrime $NC < 2$ ir $RMSEA < 0.05$.
    \item \textbf{Kuris poveikis stipriausias?}
    Žiūrėti į \texttt{Std.all} stulpelį (standartizuotus koeficientus)! Didesnis absoliutus dydis rodo didesnę santykinę įtaką.
\end{enumerate}

\end{multicols*}
\end{document}
